\section{Method}
\label{sec:method}

% Pipeline schematic (TikZ). Wrapped in \resizebox so it always fits the
% text width in both IEEEtran and the CVF/wacv two-column layout.
\begin{figure*}[t]
\centering
\resizebox{\textwidth}{!}{%
\begin{tikzpicture}[
  font=\footnotesize,
  >=Stealth,
  box/.style={draw, rounded corners, minimum height=9mm, align=center, inner sep=3pt},
  data/.style={box, fill=gray!10},
  frozen/.style={box, fill=blue!8},
  train/.style={box, fill=violet!12},
  loss/.style={box, fill=red!10}
]
% camera branch (top)
\node[data]   (img)   at (0,1.3)   {Camera\\image};
\node[frozen] (dino)  at (2.3,1.3)  {DINOv2\\ViT-B/14\\(frozen)};
\node[data]   (feat)  at (4.7,1.3)  {patch feat.\\$16{\times}16{\times}768$};
% lidar branch (bottom)
\node[data]   (lidar) at (0,-1.3)  {LiDAR\\points};
\node[data]   (vox)   at (2.3,-1.3) {voxelize\\$128^2{\times}32$};
% fused spine
\node[train]  (fuse)  at (7.2,0)   {voxel input\\$\mathbf{O}\oplus\mathbf{F}^v$};
\node[train]  (mask)  at (9.3,0)   {range-aware\\mask $\rho{=}0.85$};
\node[train]  (enc)   at (11.3,0)  {3D CNN\\encoder};
\node[train]  (bott)  at (13.0,0)  {bottleneck\\$256{\times}32^2{\times}8$};
\node[train]  (dec)   at (14.9,0)  {3D CNN\\decoder};
% heads + losses
\node[box,fill=violet!12] (occh) at (16.9,0.95) {head:\\occ.\,$\hat{\mathbf{O}}$};
\node[box,fill=violet!12] (fth)  at (16.9,-0.95) {head:\\feat.\,$\hat{\mathbf{F}}^v$};
\node[loss] (locc) at (18.9,0.95) {$\mathcal{L}_{\text{occ}}$\\BCE};
\node[loss] (lfeat) at (18.9,-0.95) {$\mathcal{L}_{\text{feat}}$\\MSE};
% arrows
\draw[->] (img)  -- (dino);
\draw[->] (dino) -- (feat);
\draw[->] (lidar)-- (vox);
\draw[->] (feat)  -- (fuse);
\draw[->] (vox)   -- (fuse);
\draw[->] (fuse) -- (mask);
\draw[->] (mask) -- (enc);
\draw[->] (enc)  -- (bott);
\draw[->] (bott) -- (dec);
\draw[->] (dec) -- (occh);
\draw[->] (dec) -- (fth);
\draw[->] (occh) -- (locc);
\draw[->] (fth)  -- (lfeat);
\end{tikzpicture}%
}
\caption{\textbf{OMU-MAE pretraining pipeline.} A paired RGB image (top)
is processed by a \emph{frozen} DINOv2 ViT-B/14 encoder; per-patch
features are projected via the camera--LiDAR calibration and
mean-aggregated per voxel. The LiDAR scan (bottom) is voxelized at
$0.4$\,m into a $128{\times}128{\times}32$ grid. The concatenated
cross-modal voxel input $\mathbf{O}\oplus\mathbf{F}^v$ is range-aware
masked ($\rho{=}0.85$), encoded by a 3D CNN, and decoded back to two
heads that reconstruct occupancy ($\mathcal{L}_{\text{occ}}$, BCE) and
DINOv2 features ($\mathcal{L}_{\text{feat}}$, MSE) at masked positions.}
\label{fig:pipeline}
\end{figure*}


Figure~\ref{fig:pipeline} gives the overall OMU-MAE pretraining pipeline.
The image branch feeds a frozen DINOv2 ViT-B/14 encoder, whose per-patch
features are projected through the camera--LiDAR calibration onto the
LiDAR voxel grid. The fused voxel input is masked, encoded by a 3D CNN,
and decoded back to two prediction heads, one for occupancy and one for
DINOv2 features, each evaluated at masked voxel positions only.

\subsection{Cross-modal voxel input}
For each scene we are given a paired RGB image
$\mathbf{I}\in\mathbb{R}^{3\times H\times W}$ and LiDAR point cloud
$\mathcal{P}=\{\mathbf{p}_i\in\mathbb{R}^4\}_{i=1}^{N}$ in the LiDAR
coordinate frame. KITTI calibration provides the projection matrices
$\mathbf{P}_2$ and $\mathbf{Tr}$ that take a LiDAR point into the
rectified left camera. We voxelize the scene over a fixed metric volume
centered on the ego-vehicle ($25.6\,\text{m}\times51.2\,\text{m}\times6.8\,\text{m}$)
at $0.4$\,m resolution, yielding a grid of size
$128\times128\times32$. Two voxel-level signals are computed:
\begin{itemize}
\item For occupancy, a binary tensor
$\mathbf{O}\in\{0,1\}^{X\times Y\times Z}$ marking voxels that contain at
least one LiDAR point.
\item For the cross-modal feature, a frozen DINOv2 ViT-B/14
encoder~\cite{dinov2} processes $\mathbf{I}$ at $224\times224$ resolution
to produce a $16\times16$ patch-feature grid
$\mathbf{F}\in\mathbb{R}^{16\times16\times768}$. For each LiDAR point we
project to image coordinates, identify the enclosing patch, and aggregate
(mean-pool) DINOv2 features into the corresponding voxel. The result is a
per-voxel feature volume $\mathbf{F}^v\in\mathbb{R}^{768\times X\times Y\times Z}$.
\end{itemize}
Voxels not visible from the camera (e.g.\ behind the vehicle, outside the
image FoV) receive zero features; this visibility mask is implicit in the
input and reproduced at decoding.

\subsection{Range-aware masking}
Following Occupancy-MAE~\cite{occmae}, voxels nearer the ego-vehicle
(where LiDAR is dense and semantic content is richest) are masked at a
higher rate. For voxel $(x,y,z)$ with metric distance $d$ from the
ego-origin we set the per-voxel mask probability
$\Pr[m_{xyz}{=}1]=\rho\cdot(1+\alpha d)^{-1}$ with global mask ratio
$\rho=0.85$ and decay $\alpha=0.5$. Both the occupancy and feature
channels are zeroed at masked locations and concatenated with a binary
mask indicator, giving an input tensor of shape
$(1+768+1)\times X\times Y\times Z$.

\subsection{Encoder, decoder architecture}
The encoder $g_\theta$ is a 4-stage 3D CNN (Conv3D $+$ GroupNorm $+$ GELU
blocks with strided downsampling) producing a bottleneck feature volume
$\mathbf{H}\in\mathbb{R}^{C\times X'\times Y'\times Z'}$ at $1/4$
resolution and $C=4\,d_{\text{feat}}=256$ channels, where
$d_{\text{feat}}=64$. The decoder mirrors the encoder via
trilinear-upsample $+$ Conv3D blocks. Two heads at the decoder output emit
(a)~a per-voxel occupancy logit and (b)~a per-voxel feature prediction in
$\mathbb{R}^{768}$.

\subsection{Dual reconstruction loss}
Let $\hat{\mathbf{O}}$ and $\hat{\mathbf{F}}^v$ denote the predicted
occupancy logits and features, and
$\mathcal{M}\in\{0,1\}^{X\times Y\times Z}$ the mask. We compute two
terms, each restricted to masked positions:
\begin{align}
\mathcal{L}_{\text{occ}} &= \frac{1}{|\mathcal{M}|}
\sum_{(x,y,z):\,m_{xyz}=1}\!\!\text{BCE}\!\left(\hat{O}_{xyz},O_{xyz}\right),\\
\mathcal{L}_{\text{feat}} &= \frac{1}{|\mathcal{M}\cap\mathcal{O}^+|}
\sum_{(x,y,z)\in\mathcal{M}\cap\mathcal{O}^+}\!\!\left\|\hat{\mathbf{F}}^v_{xyz}-\mathbf{F}^v_{xyz}\right\|_2^2,
\end{align}
where $\mathcal{O}^+=\{(x,y,z):O_{xyz}=1\}$ restricts the feature
reconstruction to \emph{occupied} voxels (predicting features for empty
space is uninformative). The combined objective is
$\mathcal{L}=\lambda_{\text{occ}}\mathcal{L}_{\text{occ}}+\lambda_{\text{feat}}\mathcal{L}_{\text{feat}}$
with $\lambda_{\text{occ}}=1.0$ and $\lambda_{\text{feat}}=0.5$. The
supervision targets are read from the un-masked voxel input (occupancy
$\mathbf{O}$ from the LiDAR voxelization, features $\mathbf{F}^v$ from the
frozen DINOv2 ViT); because the target encoder is frozen it cannot
co-adapt to the predictor, which rules out the trivial constant-output
collapse mode of joint-target prediction.
